\message{ !name(main.tex)}% 20241108
% Giovanni Ramirez
% ramirez@ecfm.usac.edu.gt
%                                                                                                                                                                                                                                                                                                                                                                                                                                                                                                                                                                                                                            
















%
% Esta es una plantilla y contiene lo mínimo que debe llevar el
% informe final del año de prácticas.  La portada es una propuesta y
% se puede ajustar según se desee.
%
% Esta plantilla se puede usar tanto en un compilador de LaTeX en su
% computadora o en algún compilador en línea como overleaf.com
%

\documentclass[titlepage,11pt]{report}

%%% PAQUETES
% esto es para hacer las cosas en español y para usar el punto como separador
% decimal
\usepackage[spanish,es-nodecimaldot]{babel}%
% esto es por la codificación con la que se escribe, revisar la codificación
% de su sistema pues podría estar en utf16 o en iso-latin-1 o iso 8859-1
\usepackage[utf8]{inputenc}%
\usepackage{float}

\usepackage{listings}
\usepackage{xcolor} % para colores

\lstdefinelanguage{Mathematica}{
  morekeywords={
    Module, Table, Plot, Integrate, Sum, If, Do, While, For, Return,
    Function, Set, SetDelayed, Rule, RuleDelayed, Block, With, True, False
  },
  sensitive=true,
  morecomment=[l](\*),
  morestring=[b]",
}

\lstset{
  language=Mathematica,
  basicstyle=\ttfamily\small,
  keywordstyle=\color{blue}\bfseries,
  commentstyle=\color{gray},
  stringstyle=\color{red},
  showstringspaces=false,
  breaklines=true,
  postbreak=\mbox{\textcolor{gray}{$\hookrightarrow$}\space}
}
% esto es para agregar gráficas y figuras
\usepackage{graphicx}%
% esto es por algunos símbolos como la ñ o las vocales tildadas
\usepackage{latexsym}%
\usepackage{amsmath, amssymb, graphics, setspace}
\usepackage{xcolor} % Para colores
\usepackage{listings} % Para mostrar código
\usepackage[many]{tcolorbox} % Para cajas con estilo
\usepackage{geometry} % Para ajustar márgenes

% Configuración de listings para Mathematica

% esto es para tener acceso a varios símbolos matemáticos
\usepackage{amsfonts, amsmath}%
% esto es para facilitar la escritura de algunos símbolos en física
\usepackage{physics}
% esto es para poder usar el sistema internacional de unidades
\usepackage[amssymb]{SIunits}%
% esto es para que la página sea tamaño carta
\usepackage[letterpaper, bottom=25mm, top=25mm]{geometry}
% esto es para los colores, la opción table es para poner colores en las
% celdas

% esto es para las tablas largas
\usepackage{longtable}
% esto es para poder escribir en varios renglones en la tabla
\usepackage{array}
\usepackage{natbib}
% esto es útil para los hipervínculos y metadata
\usepackage[bookmarksnumbered, breaklinks={true},%
pdfauthor={Bilbo Baggins (bilbo@theshire.middleearth)},%
pdftitle={El Hobbit o la historia de una ida y una vuelta},%
pdfsubject={Describir las aventuras de Bilbo},%
pdfkeywords={condensed matter materials \& applied physics, quasiparticles \&
  collective excitations, optical phonons}]{hyperref}


\begin{document}

\message{ !name(main.tex) !offset(-3) }

% % % % % % % % % % % % % % % % % % % % % % % % % % % % % % % % % % %
% aquí se define la página de la portada, las distancias y tamaños son
% relativos al tamaño de la fuente que se define al inicio en documentclass
\begin{titlepage}
  % esto es para que no ponga número a la página
  \pagestyle{empty}
  % esto agrega el título a los bookmarks
  \pdfbookmark[0]{Title}{title}
  \begin{center}
    % esto es para los logos
    \begin{minipage}{\textwidth}
      \begin{minipage}[c]{0.45\textwidth}
        \centering \includegraphics[width=.8\textwidth]{usacAzul}
      \end{minipage}%
      \hfill
      \begin{minipage}[c]{0.45\textwidth}
        \vspace{-1em}\includegraphics[width=.85\textwidth]{ecfmAzul}
      \end{minipage}
    \end{minipage}

    % esto pone los nombres de las instituciones
    \vspace{2em}
    \begin{minipage}[t]{.95\textwidth}
      \centering\large
      \textsc{Universidad de San Carlos de Guatemala\\
        Escuela de Ciencias Físicas y Matemáticas}
    \end{minipage}

    % esto pone la línea horizontal
    \vspace{1em}\noindent\rule{\textwidth}{1px}\vfill
    % esto pone el título
    \begin{minipage}[t]{.95\textwidth}
      \centering\Huge
      \textsc{Desarrollo de un modelo predictivo de precipitación para la región Bocacosta de Guatemala mediante redes neuronales LSTM}
    \end{minipage}

    % esto pone los datos 
    \vfill
    \begin{minipage}[t]{.75\textwidth}
      \centering
      Informe final de ejercicio profesional supervisado elaborado por:\\%
      \textbf{Miguel Alexander Avila Chacón}\\
      presentado ante el Departamento de Física
    \end{minipage}

    % esto pone más datos
    \vfill
    \begin{minipage}[t]{.95\textwidth}
      \centering
      Supervisado por \\%
      Ing. Klaus Hermann Aldo Johannes Zúñiga Kaufmann \\ Dr. Giovanni Ramírez
    \end{minipage}

    \vfill\noindent\rule{\textwidth}{1px}
    \begin{minipage}[t]{.95\textwidth}
      \centering
      Guatemala, junio de 2026
    \end{minipage}
  \end{center}
\end{titlepage}
% % % % % % % % % % % % % % % % % % % % % % % % % % % % % % % % % % %


% esto hace que el resumen no tenga número de página y que la introducción
% empiece con el número de página 1
\pagestyle{empty}
\setcounter{page}{0}

% % % % % % % % % % % % % % % % % % % % % % % % % % % % % % % % % % %
% en esta página se hace un resumen del trabajo, se recomienda que no exceda
% las 500 palabras.
%

\section*{Resumen}
\label{sec:resumen}

% ver metadatos
\vspace{2em}
\textbf{Palabras clave:} parton distribution functions \& phenomenology \& high energy physics.

% esto es para terminar la página del resumen y pasar a la nueva página con el
% número uno
\newpage
\pagestyle{plain}

\chapter{Introducción}
\label{sec:introduccion}

El ser humano se caracteriza por la búsqueda constante de herramientas que mejoren sus condiciones de vida. En los últimos años, estos esfuerzos se han orientado hacia la construcción de máquinas con algoritmos capaces de resolver, de forma automática y rápida, operaciones que resultan tediosas para una persona. Sin embargo, estas máquinas, basadas en algoritmos específicos, presentan una limitación importante: en ocasiones el problema que se busca resolver no requiere un tratamiento algorítmico sino el aprendizaje de patrones a través de la experiencia. 

El ser humano es capaz de resolver múltiples problemas que no pueden expresarse mediante un algoritmo gracias a su experiencia y a su capacidad para memorizar y asociar hechos. Así, parece lógico que una forma de aproximarse a este tipo de problemas consista en construir sistemas capaces de reproducir dicha característica humana. En este sentido, las redes neuronales constituyen un intento simplificado de simular el funcionamiento del cerebro humano en el tratamiento de la información, cuya unidad básica de procesamiento se inspira en la célula fundamental del sistema nervioso: la neurona.

Una característica esencial del cerebro humano es su capacidad de aprendizaje, entendida como la habilidad de resolver problemas que inicialmente no podían abordarse, mediante la adquisición de nueva información sobre ellos \citep{matich2001redes, salas2004redes, tablada2009redes, izaurieta2000redes}.

Por lo tanto, las redes neuronales se componen de unidades de procesamiento que intercambian datos o información. Se utilizan para reconocer patrones en conjuntos de datos, como imágenes, manuscritos o secuencias temporales, con el objetivo de aprender de ellos y mejorar su desempeño \citep{matich2001redes}.

Debido a sus fundamentos, las redes neuronales artificiales presentan un gran número de características semejantes a las del cerebro, y por ello una gran cantidad de ventajas por sobre los programas basados en algoritmos \cite{McCullochPitts1943}. Por ejemplo, son capaces de aprender a través de la experiencia, de generalizar de una gran serie de casos específicos a nuevos casos, de abstraer características esenciales de bancos de información, desechando la información irrelevante, etc. 

\chapter{Objetivos}
\label{sec:objetivos}
\section{General}
\label{sec:general}

Desarrollar un modelo predictivo de precipitación para la región de Bocacosta en Guatemala utilizando redes neuronales LSTM. Con datos históricos con el propósito de complementar los métodos tradicionales de pronóstico climático. 

\section{Específicos}
\label{sec:especificos}

\begin{enumerate}
    \item Analizar series de tiempo de variables climáticas, como precipitación y temperatura, utilizando estadística descriptiva. 
    \item Diseñar e implementar un modelo LSTM utilizando datos históricos del INSIVUMEH.
    \item Evaluar el desempeño del modelo mediante métricas de validación, como el análisis de correlación.
    \item Comparar los resultados obtenidos con métodos tradicionales.
\end{enumerate}

\chapter{Marco teórico}
\label{sec:marco}


\section{Dinámica de la Atmósfera}

\subsection{Conceptos Básicos del Sistema Climático}
El sistema climático terrestre es un sistema acoplado no lineal compuesto por la atmósfera, la hidrósfera, la criósfera, la litósfera y la biósfera. Su dinámica está gobernada por las leyes de conservación de la masa, el momento y la energía, así como por la ecuación de estado de los gases ideales. La fuente primaria de energía es la radiación solar de onda corta, la cual es distribuida de manera desigual sobre la superficie terrestre debido a la esfericidad del planeta y la inclinación del eje de rotación. Este desequilibrio radiativo genera gradientes de temperatura que impulsan la circulación general de la atmósfera y los océanos, cuyo objetivo termodinámico es transportar calor desde las regiones ecuatoriales hacia los polos.

\subsection{Variables Atmosférias Relevantes}
Para caracterizar el estado de la atmósfera y su evolución, se utilizan variables de estado fundamentales. La \textbf{presión atmosférica} (\(P\)), definida como la fuerza por unidad de área ejercida por el peso de la columna de aire, decrece exponencialmente con la altura siguiendo la ecuación hipsométrica. La \textbf{temperatura} (\(T\)) es una medida de la energía cinética media de las moléculas del aire. El \textbf{contenido de vapor de agua} se cuantifica mediante la humedad específica (\(q\)), que representa la masa de vapor de agua por unidad de masa de aire húmedo, o la humedad relativa (\(HR\)), que indica el grado de saturación del aire respecto a la presión de vapor de saturación. El \textbf{viento}, un campo vectorial \(\vec{V} = (u, v, w)\), describe el movimiento del aire en las direcciones zonal, meridional y vertical, respectivamente.

\subsection{Procesos Termodinámicos}
Los procesos adiabáticos son centrales en la formación de nubes y precipitación. Cuando una parcela de aire asciende en la atmósfera, encuentra una presión ambiental menor y se expande. Si el ascenso es lo suficientemente rápido para que el intercambio de calor con el entorno sea despreciable (proceso adiabático), la parcela se enfría a la tasa adiabática seca (\(\Gamma_d \approx 9.8 \text{ K km}^{-1}\)). Una vez que la parcela alcanza la saturación, la condensación libera calor latente, reduciendo la tasa de enfriamiento a la tasa pseudo-adiabática húmeda (\(\Gamma_s\)), cuyo valor varía con la temperatura y la presión, siendo típicamente de \(4-6 \text{ K km}^{-1}\) en la troposfera tropical. Este calor latente es una fuente crucial de energía para la convección profunda.

\section{Mecanismos de Formación de la Precipitación}

\subsection{Saturación, Condensación y Formación de Nubes}
La precipitación es el resultado final de una cadena de procesos microfísicos que comienzan con la saturación del aire. La \textbf{ecuación de Clausius-Clapeyron} establece la relación termodinámica fundamental entre la presión de vapor de saturación (\(e_s\)) y la temperatura:
\begin{equation}
    \frac{de_s}{dT} = \frac{L_v e_s}{R_v T^2}
\end{equation}
donde \(L_v\) es el calor latente de vaporización y \(R_v\) es la constante específica del vapor de agua. Esta ecuación implica que la capacidad del aire para contener vapor de agua aumenta exponencialmente con la temperatura. En un clima más cálido, la atmósfera puede retener más humedad, lo que intensifica los eventos extremos de precipitación si los mecanismos de ascenso están presentes. La condensación del vapor sobre núcleos de condensación de nubes (CCN) da lugar a las gotitas de nube, cuyo tamaño es insuficiente para precipitar por gravedad.

\subsection{Procesos de Colisión-Coalescencia y Bergeron-Findeisen}
La formación de gotas de lluvia requiere el crecimiento de las gotitas de nube hasta alcanzar un tamaño crítico. En nubes cálidas (temperatura \(> 0^\circ\)C), predomina el proceso de \textbf{colisión-coalescencia}, donde gotas de diferente tamaño y velocidad terminal colisionan y se fusionan. En nubes frías o mixtas, opera el \textbf{proceso de Bergeron-Findeisen}, que explota la diferencia de presión de vapor de saturación sobre el agua líquida y el hielo. En un ambiente mixto saturado respecto al agua, existe sobresaturación respecto al hielo, lo que induce la deposición de vapor sobre los cristales de hielo a expensas de las gotitas de agua subfundida.

\subsection{Ascenso de Masas de Aire}
Para que se produzca precipitación significativa, es indispensable el ascenso de masas de aire húmedo que enfríen la parcela hasta la saturación. Los mecanismos de ascenso principales son:
\begin{itemize}
    \item \textbf{Convección libre:} Calentamiento diferencial de la superficie que genera burbujas de aire cálido menos denso que ascienden.
    \item \textbf{Ascenso orográfico:} El viento horizontal es forzado a remontar una barrera montañosa, enfriándose adiabáticamente. Este mecanismo es crítico en la Bocacosta guatemalteca.
    \item \textbf{Convergencia en niveles bajos:} La fricción superficial o la dinámica de ondas del este generan zonas de acumulación de masa que fuerzan el movimiento vertical ascendente.
    \item \textbf{Frentes:} Superficies de discontinuidad entre masas de aire de distinta densidad.
\end{itemize}

\section{Variables Físicas Relevantes para la Predicción de Precipitación}

La selección de predictores (características de entrada) para un modelo de precipitación debe basarse en su conexión causal con los procesos de condensación y ascenso.

\subsection{Temperatura Superficial del Mar (SST)}
La SST es la variable oceánica más influyente en el clima tropical. Anomalías cálidas en la SST incrementan el flujo de calor latente y sensible hacia la atmósfera, reduciendo la estabilidad estática e incrementando la humedad en la capa límite. En el contexto del Pacífico Oriental, la SST modula la intensidad de la Zona de Convergencia Intertropical (ZCIT) y la convección asociada a ondas del este.

\subsection{Humedad Relativa y Específica}
La humedad específica (\(q\)) en niveles bajos y medios (850-700 hPa) es un indicador directo del combustible disponible para la convección. La humedad relativa (\(HR\)) proporciona información sobre qué tan cerca está la columna atmosférica de la saturación, lo cual es crucial para la activación de la convección profunda.

\subsection{Viento (Componentes Zonal y Meridional)}
El viento zonal (\(u\)) en niveles bajos (925-850 hPa) es fundamental para el transporte de humedad desde las fuentes oceánicas (el Jet de Bajo Nivel del Caribe o el flujo del Pacífico). La componente meridional (\(v\)) está asociada a la convergencia/divergencia y a la posición de la ZCIT. La cizalladura vertical del viento puede tanto inhibir como organizar los sistemas convectivos.

\subsection{Geopotencial}
La altura geopotencial (\(Z\)) en niveles de 500 hPa permite identificar la presencia de vaguadas y dorsales en la tropósfera media. Una vaguada profunda sobre Centroamérica se asocia con ascenso de gran escala y enfriamiento en la tropósfera superior, lo que desestabiliza la columna atmosférica y favorece tormentas.

\subsection{Teleconexiones Climáticas (ENSO, AMO, TNA)}
El sistema climático presenta modos de variabilidad de baja frecuencia que modulan las condiciones medias de precipitación en Guatemala.
\begin{itemize}
    \item \textbf{ENSO (El Niño-Southern Oscillation):} La fase cálida (El Niño) se asocia típicamente con déficit de precipitación en la vertiente del Pacífico guatemalteco debido al debilitamiento de los vientos alisios y al desplazamiento de la convección profunda hacia el Pacífico Central. La fase fría (La Niña) se correlaciona con excesos de lluvia.
    \item \textbf{AMO (Atlantic Multidecadal Oscillation):} Afecta la SST del Atlántico Norte y modula la fuerza del transporte de humedad desde el Caribe hacia el Pacífico a través del istmo.
    \item \textbf{TNA (Tropical North Atlantic Index):} Anomalías cálidas en el Atlántico tropical intensifican la corriente de humedad hacia Centroamérica, compitiendo con la influencia del Pacífico.
\end{itemize}

\section{Series de Tiempo en Climatología}

\subsection{Definición y Componentes}
Una serie de tiempo climática es una secuencia de observaciones de una variable atmosférica ordenadas cronológicamente y generalmente equiespaciadas. Matemáticamente, se define como \(\{X_t : t \in T\}\), donde \(T\) es el índice temporal. Estas series suelen descomponerse en tres componentes aditivos o multiplicativos:
\begin{equation}
    X_t = T_t + S_t + R_t
\end{equation}
donde \(T_t\) es la \textbf{tendencia} a largo plazo (cambio climático o variabilidad decadal), \(S_t\) es la \textbf{estacionalidad} periódica (ciclo anual marcado por el monzón centroamericano), y \(R_t\) es el residuo o \textbf{variabilidad} de alta frecuencia (sinóptica e intradiurna).

\subsection{Autocorrelación}
La función de autocorrelación (ACF) mide la dependencia lineal entre los valores de la serie en el tiempo \(t\) y en rezagos \(t-k\). En climatología, la precipitación exhibe una \textbf{persistencia} significativa; la lluvia de hoy está correlacionada con la lluvia de ayer, y más aún, con la humedad del suelo o los modos de variabilidad oceánica de semanas previas. Esta propiedad justifica el uso de modelos con memoria temporal como las redes recurrentes.

\section{Relación entre Procesos Físicos y Aprendizaje Automático}

\subsection{Naturaleza No Lineal del Sistema Climático}
La atmósfera es un fluido geofísico turbulento gobernado por las ecuaciones de Navier-Stokes no lineales. Los procesos de retroalimentación, como el acoplamiento entre la radiación de onda larga y la nubosidad, o la liberación de calor latente en función de la convergencia de humedad, generan una respuesta altamente no lineal del sistema. Esta no linealidad implica que pequeños errores en las condiciones iniciales crecen exponencialmente (caos determinista), limitando la predictibilidad de los modelos dinámicos más allá de unos pocos días para el detalle convectivo.

\subsection{Limitaciones de Modelos Físicos Tradicionales}
Los Modelos de Circulación General (GCM) y los Modelos de Predicción Numérica del Tiempo (NWP) resuelven las ecuaciones primitivas en una malla discreta. Sin embargo, enfrentan dificultades inherentes en regiones de terreno complejo como la Bocacosta, donde la resolución espacial es insuficiente para resolver la orografía escarpada y los procesos de convección sub-malla deben ser \textbf{parametrizados} (aproximados empíricamente). Estas parametrizaciones introducen sesgos sistemáticos y errores en el pronóstico de precipitación orográfica.

\subsection{Justificación del Uso de Machine Learning}
El Aprendizaje Automático (ML), y en particular las redes neuronales profundas, ofrecen un enfoque complementario puramente basado en datos. En lugar de parametrizar la microfísica, una red neuronal puede aprender empíricamente la función de mapeo \(f: \mathbb{R}^n \to \mathbb{R}\) que relaciona un conjunto de predictores físicos de gran escala (resueltos por reanálisis) con la precipitación local observada. Este enfoque permite capturar relaciones estadísticas no lineales y dependencias temporales complejas que son difíciles de formular explícitamente en un modelo físico.

\subsection{Interpretabilidad Física de los Predictores}
Si bien las redes neuronales son a menudo criticadas como ``cajas negras'', la selección cuidadosa de predictores basada en la dinámica atmosférica (Sección 3) permite una \textbf{interpretabilidad física a priori}. Sabemos que la red está utilizando información sobre la SST del Pacífico, la humedad transportada por el viento zonal y la estabilidad atmosférica para inferir la probabilidad e intensidad de la precipitación, en concordancia con la teoría de la convección.

\section{Redes Neuronales Artificiales}

\subsection{Definición Formal y Modelo Matemático}
Una red neuronal artificial es un modelo computacional inspirado en el sistema nervioso biológico, compuesto por unidades de procesamiento interconectadas llamadas neuronas. Formalmente, una neurona artificial recibe un vector de entrada \(\mathbf{x} = (x_1, x_2, ..., x_n)\), calcula una suma ponderada de estas entradas y aplica una función de activación no lineal \(g(\cdot)\):
\begin{equation}
    y = g\left( \sum_{i=1}^{n} w_i x_i + b \right) = g(\mathbf{w}^T \mathbf{x} + b)
\end{equation}
donde \(\mathbf{w}\) son los pesos sinápticos (parámetros ajustables que determinan la importancia de cada entrada) y \(b\) es el sesgo o \textit{bias}.

\subsection{Funciones de Activación}
La no linealidad \(g(\cdot)\) es esencial para que la red pueda aproximar funciones complejas. Las funciones más comunes en modelos LSTM son:
\begin{itemize}
    \item \textbf{Sigmoide (\(\sigma\)):} \(\sigma(x) = \frac{1}{1+e^{-x}}\). Comprime los valores al rango \((0, 1)\), utilizada en puertas de control de flujo de información.
    \item \textbf{Tangente Hiperbólica (\(\tanh\)):} \(\tanh(x) = \frac{e^x - e^{-x}}{e^x + e^{-x}}\). Comprime valores al rango \((-1, 1)\), centrada en cero, utilizada para regular el estado de la celda.
    \item \textbf{ReLU (Rectified Linear Unit):} \(ReLU(x) = \max(0, x)\). Popular en capas ocultas densas por su eficiencia computacional y mitigación del desvanecimiento del gradiente.
\end{itemize}

\section{Redes Neuronales Recurrentes (RNN)}

\subsection{Concepto de Memoria Temporal}
Las redes neuronales \textit{feedforward} estándar asumen independencia entre muestras, lo cual es inadecuado para series de tiempo donde \(X_t\) depende de \(X_{t-1}\). Las Redes Neuronales Recurrentes (RNN) introducen conexiones de retroalimentación que forman un ciclo dirigido, permitiendo que la información persista. El estado oculto \(\mathbf{h}_t\) actúa como una ``memoria'' que resume la historia pasada de la secuencia:
\begin{equation}
    \mathbf{h}_t = g(\mathbf{W}_{xh} \mathbf{x}_t + \mathbf{W}_{hh} \mathbf{h}_{t-1} + \mathbf{b}_h)
\end{equation}
donde \(\mathbf{W}_{hh}\) es la matriz de pesos recurrente que propaga el estado anterior.

\subsection{Problema del Gradiente Evanescente}
El entrenamiento de RNNs mediante retropropagación a través del tiempo (BPTT) implica el cálculo de gradientes a través de múltiples pasos temporales. Debido a multiplicaciones repetidas de la matriz \(\mathbf{W}_{hh}\) y las derivadas de las funciones de activación, los gradientes tienden a \textbf{desvanecerse} (tender a cero) o \textbf{explotar} (tender a infinito). En el contexto climático, donde la precipitación actual puede estar influenciada por la SST de hace meses (teleconexiones ENSO), la RNN simple es incapaz de capturar estas dependencias de \textbf{largo plazo} debido al gradiente evanescente.

\section{Redes LSTM (Long Short-Term Memory)}

\subsection{Justificación de su Uso}
La arquitectura LSTM fue diseñada específicamente para superar el problema del gradiente evanescente y aprender dependencias temporales de largo alcance. Dado que los fenómenos atmosféricos operan en múltiples escalas temporales (desde la turbulencia horaria hasta la oscilación interanual del ENSO), la LSTM es la elección óptima para modelar la precipitación acumulada. Su capacidad para recordar y olvidar información selectivamente permite a la red mantener un estado que refleje la condición climática de fondo (e.g., ``estamos en fase La Niña'') mientras reacciona a perturbaciones sinópticas de corta duración.

\subsection{Arquitectura Interna y Ecuaciones Fundamentales}
La celda LSTM reemplaza la neurona recurrente simple por un módulo que contiene una \textbf{celda de estado} (\(C_t\)) y tres compuertas que regulan el flujo de información.

\begin{enumerate}
    \item \textbf{Puerta de Olvido (\(f_t\)):} Decide qué información eliminar del estado de la celda previo \(C_{t-1}\). Observa \(\mathbf{h}_{t-1}\) y \(\mathbf{x}_t\), y emite un valor entre 0 (olvidar) y 1 (retener).
    \begin{equation}
        f_t = \sigma(\mathbf{W}_f \cdot [\mathbf{h}_{t-1}, \mathbf{x}_t] + \mathbf{b}_f)
    \end{equation}

    \begin{equation*}
      f^2=2x+2
    \end{equation*}
    
    \item \textbf{Puerta de Entrada (\(i_t\)) y Candidato de Celda (\(\tilde{C}_t\)):} Decide qué nueva información almacenar en la celda. La puerta de entrada selecciona los valores a actualizar, mientras que una capa \(\tanh\) genera un vector de nuevos valores candidatos.
    \begin{equation}
        i_t = \sigma(\mathbf{W}_i \cdot [\mathbf{h}_{t-1}, \mathbf{x}_t] + \mathbf{b}_i)
    \end{equation}
    \begin{equation}
        \tilde{C}_t = \tanh(\mathbf{W}_C \cdot [\mathbf{h}_{t-1}, \mathbf{x}_t] + \mathbf{b}_C)
    \end{equation}
    
    \item \textbf{Actualización del Estado de Celda (\(C_t\)):} El estado anterior se actualiza combinando la acción del olvido y la adición de la nueva información.
    \begin{equation}
        C_t = f_t * C_{t-1} + i_t * \tilde{C}_t
    \end{equation}
    El flujo lineal de gradiente a través de la operación \(+\) y \(*\) (producto elemento a elemento) es lo que evita el desvanecimiento del gradiente a lo largo de decenas de pasos temporales.

    \item \textbf{Puerta de Salida (\(o_t\)) y Estado Oculto (\(h_t\)):} La salida se basa en el estado de la celda, pero filtrado por una compuerta de salida que decide qué partes del estado actual son relevantes para el siguiente paso o la predicción final.
    \begin{equation}
        o_t = \sigma(\mathbf{W}_o \cdot [\mathbf{h}_{t-1}, \mathbf{x}_t] + \mathbf{b}_o)
    \end{equation}
    \begin{equation}
        \mathbf{h}_t = o_t * \tanh(C_t)
    \end{equation}
\end{enumerate}

\section{Preprocesamiento de Datos Climáticos}

\subsection{Normalización y Estandarización}
Las variables climáticas poseen escalas y unidades físicas dispares (e.g., geopotencial en \(m^2 s^{-2}\), humedad relativa en \%). Para que el optimizador por gradiente converja eficientemente y las funciones de activación operen en su rango sensible, es imperativo escalar los datos. La \textbf{estandarización} (o puntuación Z) se utiliza comúnmente en climatología para preservar las anomalías relativas:
\begin{equation}
    X_{std} = \frac{X - \mu}{\sigma}
\end{equation}
donde \(\mu\) es la media y \(\sigma\) la desviación estándar del período de entrenamiento.

\subsection{Manejo de Datos Faltantes}
Los conjuntos de datos observacionales (pluviómetros) o satelitales contienen lagunas. Se aplican técnicas de imputación como la interpolación lineal temporal o la sustitución por la mediana climatológica del mes correspondiente. Para mantener la consistencia física, es preferible eliminar segmentos cortos de reanálisis en lugar de interpolar valores espurios de viento.

\subsection{Construcción de Ventanas Temporales (Lags)}
Para alimentar la red LSTM, la serie de tiempo continua debe ser segmentada en secuencias de longitud fija \(L\) (ventana de entrada). Si se predice la precipitación en el tiempo \(t\), la entrada es la matriz \(\mathbf{X} \in \mathbb{R}^{L \times N}\) donde \(N\) es el número de predictores. La elección de \(L\) es crítica: debe ser lo suficientemente larga para capturar la evolución de una onda del este (3-5 días) y la memoria de la humedad del suelo, pero no tan larga como para diluir el gradiente en el entrenamiento.

\subsection{División de Datos}
Para evaluar la capacidad de generalización del modelo, el conjunto de datos se divide en:
\begin{itemize}
    \item \textbf{Entrenamiento:} Aproximadamente 70-80\% de los datos (e.g., años iniciales). Usado para ajustar los pesos \(\mathbf{W}\) y \(\mathbf{b}\) mediante retropropagación.
    \item \textbf{Validación:} 10-15\% de los datos. Usado para monitorear el rendimiento fuera de muestra y detener el entrenamiento (\textit{early stopping}) antes del sobreajuste.
    \item \textbf{Prueba:} 10-15\% final. Conjunto de datos \textbf{no vistos} durante el desarrollo. Proporciona una estimación insesgada del rendimiento operativo del modelo final.
\end{itemize}

\section{Evaluación de Modelos Predictivos}

La bondad del ajuste y la precisión del pronóstico se cuantifican mediante métricas de error y correlación.

\subsection{RMSE (Raíz del Error Cuadrático Medio)}
Mide la desviación estándar de los residuos (errores de predicción). Penaliza fuertemente los errores grandes debido al término cuadrático. Es útil para cuantificar la magnitud media del error en las mismas unidades que la variable observada (mm/día).
\begin{equation}
    RMSE = \sqrt{\frac{1}{n} \sum_{i=1}^{n} (y_i - \hat{y}_i)^2}
\end{equation}

\subsection{MAE (Error Absoluto Medio)}
Mide la magnitud promedio de los errores en un conjunto de predicciones, sin considerar la dirección. Es menos sensible a valores atípicos extremos que el RMSE.
\begin{equation}
    MAE = \frac{1}{n} \sum_{i=1}^{n} |y_i - \hat{y}_i|
\end{equation}

\subsection{Correlación de Pearson (\(r\))}
Evalúa la relación lineal entre los valores predichos (\(\hat{y}\)) y los observados (\(y\)). Un valor cercano a 1 indica que el modelo captura correctamente la fase y la tendencia de la variabilidad de la precipitación, aunque pudiera tener un sesgo sistemático en la magnitud.
\begin{equation}
    r = \frac{\sum (y_i - \bar{y})(\hat{y}_i - \bar{\hat{y}})}{\sqrt{\sum (y_i - \bar{y})^2 \sum (\hat{y}_i - \bar{\hat{y}})^2}}
\end{equation}

\section{Problemas en el Entrenamiento de Modelos}

\subsection{Overfitting (Sobreajuste)}
Ocurre cuando el modelo aprende no solo la señal física subyacente, sino también el \textit{ruido} estadístico y las idiosincrasias del conjunto de entrenamiento. Un modelo sobreajustado exhibe un error muy bajo en entrenamiento pero un error alto en validación/prueba. En el contexto de LSTM, el sobreajuste se manifiesta cuando la red memoriza secuencias específicas de años pasados en lugar de aprender las relaciones atmósfera-lluvia generalizables.

\subsection{Underfitting (Subajuste)}
Sucede cuando el modelo es demasiado simple (pocas neuronas o capas) para capturar la complejidad no lineal del sistema climático. El error es alto tanto en entrenamiento como en validación, indicando que el modelo no ha logrado aprender la función de mapeo verdadera.

\subsection{Técnicas de Regularización}
Para combatir el sobreajuste en redes LSTM profundas se emplean:
\begin{itemize}
    \item \textbf{Dropout:} Durante el entrenamiento, se ``apaga'' aleatoriamente un porcentaje \(p\) de las neuronas en cada paso (ignorando su salida y gradiente). Esto fuerza a la red a no depender de una ruta de activación específica, mejorando la robustez y actuando como un ensamble de sub-redes.
    \item \textbf{Early Stopping:} El entrenamiento se detiene cuando el error de validación deja de disminuir y comienza a aumentar, incluso si el error de entrenamiento sigue bajando. Es una de las técnicas más efectivas en series temporales climáticas.
\end{itemize}

\section{Contexto Climático de Guatemala y Región Bocacosta}

\subsection{Estacionalidad de la Precipitación}
La región Bocacosta se ubica en la vertiente sur de la cadena volcánica de Guatemala, descendiendo hacia la llanura costera del Pacífico. El régimen de precipitación es marcadamente estacional, definido por el monzón norteaméricano y la migración latitudinal de la zona de convergencia intertropical. La temporada lluviosa se extiende típicamente de mayo a octubre, presentando un característico \textbf{doble máximo}: un primer pico en junio asociado al inicio de la época lluviosa y al paso de ondas del este, una \textbf{canícula} en julio-agosto debido a la intensificación de los vientos alisios y subsidencia, y un segundo máximo en septiembre-octubre coincidiendo con el pico de actividad de la zona de convergencia intertropical en el Pacífico Oriental y la temporada de ciclones tropicales.

\subsection{Influencia del Pacífico y Rol de la Orografía}
La Bocacosta está directamente expuesta al flujo de humedad del Océano Pacífico tropical nororiental. La interacción entre el flujo de bajo nivel del suroeste y la pronunciada orografía de la Sierra Madre guatemalteca es el principal mecanismo forzante de la precipitación en la zona. El ascenso forzado (efecto Föehn) condensa la abundante humedad oceánica, generando núcleos de precipitación convectiva\begin{comment}y estratiforme que superan los 4000 mm anuales en las laderas de barlovento\end{comment}. Pequeñas variaciones en la dirección del viento o en la temperatura superficial del Pacífico pueden traducirse en grandes anomalías de precipitación local, exacerbando el riesgo de deslizamientos e inundaciones.

\subsection{Importancia de la Predicción en la Región}
La Bocacosta es una de las zonas con más producción agrícola del país (café, caña de azúcar, hule), pero también una de las más vulnerables a eventos hidrometeorológicos extremos. La predicción precisa de la precipitación a corto y mediano plazo es un insumo fundamental para:
\begin{itemize}
    \item Sistemas de Alerta Temprana (SAT) ante deslizamientos y crecidas de ríos.
    \item Planificación de siembra y cosecha en el sector agrícola.
    \item Gestión de recursos hídricos para generación hidroeléctrica en la vertiente del Pacífico.
\end{itemize}
El desarrollo de un modelo LSTM para la precipitación local busca complementar las limitaciones de los modelos numéricos globales, proporcionando una herramienta basada en predictores físicos de gran escala.





\chapter{Metodología}
\label{sec:metodologia}



\chapter{Resultados}
\label{sec:resultados}



\chapter{Conclusiones}
\label{sec:conclusiones}


\chapter{Recomendaciones}
\label{sec:recomendaciones}

% adjunto con el nombre referencias.bib.  Hay varios estilos, los básicos son:
% abbrv, acm, alpha, apalike, ieeetr, plain, siam, unsrt
\nocite{*}
\bibliographystyle{unsrt}% otras opciones {plain}
\bibliography{biblio}
\end{document}


%%% Local Variables:
%%% mode: latex
%%% TeX-master: t
%%% End:


\message{ !name(main.tex) !offset(-488) }
